\section{Regular local rings}
\label{section-regular}

\noindent
Regular local rings are defined in Definition \ref{definition-regular-local}.
It is not that easy to show that all prime localizations of a regular local
ring are regular. In fact, quite a bit of the material developed so far is
geared towards a proof of this fact. See
Proposition \ref{proposition-finite-gl-dim-regular}, and
trace back the references.

\begin{lemma}
\label{lemma-regular-graded}
Let $(R, \mathfrak m, \kappa)$ be a regular local ring of dimension $d$.
The graded ring $\bigoplus \mathfrak m^n / \mathfrak m^{n + 1}$
is isomorphic to the graded polynomial algebra
$\kappa[X_1, \ldots, X_d]$.
\end{lemma}

\begin{proof}
Let $x_1, \ldots, x_d$ be a minimal set of generators
for the maximal ideal $\mathfrak m$, see
Definition \ref{definition-regular-local}.
There is a surjection $\kappa[X_1, \ldots, X_d]
\to \bigoplus \mathfrak m^n/\mathfrak m^{n + 1}$,
which maps $X_i$ to the class of $x_i$ in $\mathfrak m/\mathfrak m^2$.
Since $d(R) = d$ by
Proposition \ref{proposition-dimension}
we know that the numerical
polynomial $n \mapsto \dim_\kappa \mathfrak m^n/\mathfrak m^{n + 1}$
has degree $d - 1$. By Lemma \ref{lemma-quotient-smaller-d} we
conclude that the surjection $\kappa[X_1, \ldots, X_d]
\to \bigoplus \mathfrak m^n/\mathfrak m^{n + 1}$ is an isomorphism.
\end{proof}

\begin{lemma}
\label{lemma-regular-domain}
Any regular local ring is a domain.
\end{lemma}

\begin{proof}
We will use that $\bigcap \mathfrak m^n = 0$
by Lemma \ref{lemma-intersect-powers-ideal-module-zero}.
Let $f, g \in R$ such that $fg = 0$.
Suppose that $f \in \mathfrak m^a$ and
$g \in \mathfrak m^b$, with $a, b$ maximal.
Since $fg = 0 \in \mathfrak m^{a + b + 1}$
we see from the result of Lemma \ref{lemma-regular-graded}
that either $f \in \mathfrak m^{a + 1}$ or
$g \in \mathfrak m^{b + 1}$. Contradiction.
\end{proof}

\begin{lemma}
\label{lemma-regular-ring-CM}
Let $R$ be a regular local ring and let
$x_1, \ldots, x_d$ be a minimal set of generators
for the maximal ideal $\mathfrak m$. Then
$x_1, \ldots, x_d$ is a regular sequence, and
each $R/(x_1, \ldots, x_c)$ is a regular local ring
of dimension $d - c$. In particular $R$ is Cohen-Macaulay.
\end{lemma}

\begin{proof}
Note that $R/x_1R$ is a Noetherian local ring of dimension $\geq d - 1$
by Lemma \ref{lemma-one-equation} with $x_2, \ldots, x_d$
generating the maximal ideal. Hence it is a regular local ring by definition.
Since $R$ is a domain by Lemma \ref{lemma-regular-domain}
$x_1$ is a nonzerodivisor.
\end{proof}

\begin{lemma}
\label{lemma-regular-quotient-regular}
Let $R$ be a regular local ring. Let $I \subset R$ be an ideal
such that $R/I$ is a regular local ring as well. Then
there exists a minimal set of generators $x_1, \ldots, x_d$
for the maximal ideal $\mathfrak m$ of $R$ such that
$I = (x_1, \ldots, x_c)$ for some $0 \leq c \leq d$.
\end{lemma}

\begin{proof}
Say $\dim(R) = d$ and $\dim(R/I) = d - c$.
Denote $\overline{\mathfrak m} = \mathfrak m/I$ the
maximal ideal of $R/I$. Let $\kappa = R/\mathfrak m$. We have
$$
\dim_\kappa((I + \mathfrak m^2)/\mathfrak m^2) =
\dim_\kappa(\mathfrak m/\mathfrak m^2)
- \dim(\overline{\mathfrak m}/\overline{\mathfrak m}^2) = d - (d - c) = c
$$
by the definition of a regular local ring. Hence we can choose
$x_1, \ldots, x_c \in I$ whose images in $\mathfrak m/\mathfrak m^2$
are linearly independent and supplement with
$x_{c + 1}, \ldots, x_d$ to get a minimal system of generators of
$\mathfrak m$. The induced map $R/(x_1, \ldots, x_c) \to R/I$ is a
surjection between regular local rings of the same dimension
(Lemma \ref{lemma-regular-ring-CM}). It follows that
the kernel is zero, i.e., $I = (x_1, \ldots, x_c)$. Namely, if not
then we would have $\dim(R/I) < \dim(R/(x_1, \ldots, x_c))$ by
Lemmas \ref{lemma-regular-domain} and \ref{lemma-one-equation}.
\end{proof}

\begin{lemma}
\label{lemma-free-mod-x}
Let $R$ be a Noetherian local ring.
Let $x \in \mathfrak m$.
Let $M$ be a finite $R$-module such that
$x$ is a nonzerodivisor on $M$ and
$M/xM$ is free over $R/xR$.
Then $M$ is free over $R$.
\end{lemma}

\begin{proof}
Let $m_1, \ldots, m_r$ be elements of $M$ which map to
a $R/xR$-basis of $M/xM$. By Nakayama's Lemma \ref{lemma-NAK}
$m_1, \ldots, m_r$ generate $M$. If $\sum a_i m_i = 0$
is a relation, then $a_i \in xR$ for all $i$. Hence
$a_i = b_i x$ for some $b_i \in R$. Hence
the kernel $K$ of $R^r \to M$ satisfies $xK = K$
and hence is zero by Nakayama's lemma.
\end{proof}

\begin{lemma}
\label{lemma-regular-mcm-free}
Let $R$ be a regular local ring.
Any maximal Cohen-Macaulay module over $R$ is free.
\end{lemma}

\begin{proof}
Let $M$ be a maximal Cohen-Macaulay module over $R$.
Let $x \in \mathfrak m$ be part of a regular sequence
generating $\mathfrak m$. Then $x$ is a nonzerodivisor
on $M$ by Proposition \ref{proposition-CM-module}, and
$M/xM$ is a maximal Cohen-Macaulay module over $R/xR$.
By induction on $\dim(R)$ we see that $M/xM$ is free.
We win by Lemma \ref{lemma-free-mod-x}.
\end{proof}

\begin{lemma}
\label{lemma-regular-mod-x}
Suppose $R$ is a Noetherian local ring.
Let $x \in \mathfrak m$ be a nonzerodivisor
such that $R/xR$ is a regular local ring. Then $R$ is a regular local ring.
More generally, if $x_1, \ldots, x_r$ is a regular sequence in $R$
such that $R/(x_1, \ldots, x_r)$ is a regular local ring, then
$R$ is a regular local ring.
\end{lemma}

\begin{proof}
This is true because $x$ together with the lifts of a system
of minimal generators of the maximal ideal of $R/xR$ will give
$\dim(R)$ generators of $\mathfrak m$.
Use Lemma \ref{lemma-one-equation}.
The last statement follows from the first and induction.
\end{proof}

\begin{lemma}
\label{lemma-colimit-regular}
Let $(R_i, \varphi_{ii'})$ be a directed system of local rings whose
transition maps are local ring maps. If each $R_i$ is a regular local ring and
$R = \colim R_i$ is Noetherian, then $R$ is a regular local ring.
\end{lemma}

\begin{proof}
Let $\mathfrak m \subset R$ be the maximal ideal; it is the colimit
of the maximal ideal $\mathfrak m_i \subset R_i$.
We prove the lemma by induction on $d = \dim \mathfrak m/\mathfrak m^2$.
If $d = 0$, then $R = R/\mathfrak m$ is a field and $R$ is a regular local ring.
If $d > 0$ pick an $x \in \mathfrak m$, $x \not \in \mathfrak m^2$.
For some $i$ we can find an $x_i \in \mathfrak m_i$ mapping to $x$.
Note that $R/xR = \colim_{i' \geq i} R_{i'}/x_iR_{i'}$ is a Noetherian
local ring. By
Lemma \ref{lemma-regular-ring-CM}
we see that $R_{i'}/x_iR_{i'}$ is a regular local ring.
Hence by induction we see
that $R/xR$ is a regular local ring. Since each $R_i$ is a domain
(Lemma \ref{lemma-regular-graded}) we see that $R$ is a domain.
Hence $x$ is a nonzerodivisor and we conclude that $R$ is
a regular local ring by Lemma \ref{lemma-regular-mod-x}.
\end{proof}





